\documentclass[11pt, letterpaper]{article}
\usepackage{graphicx} % Required for inserting images
\usepackage{biblatex} % For my references
\addbibresource{references.bib}
\title{DOI Working Title: The Effect of Network Organization on Diffusion of Innovation in the Healthcare Environment}
\author{\textit{Sydney Corum, Anagh Mishra, Seth Tumlin}}
\date{\today}

\begin{document}

\maketitle
\begin{abstract}
    
\end{abstract}
\section{Introduction}
 The basis of this topic is found in the Diffusion of Innovations Theory by Everett M. Rogers The four pillars of the diffusion of innovation are the innovation or idea itself, the communication channels in which it can spread, the time in which the idea spreads and the time in the decision process, and the complex social network that the innovation is birthed into.\cite{rogers_diffusion_1983} Each of these pillars have their own characteristics that make modeling the diffusion of an innovation difficult.  

Understanding how new innovations, whether products or technological concepts, spread is critical, as diffusion models can help explain why some ideas successfully reach the target population, whereas others remain stagnant. These innovations arise across diverse domains, including agriculture, consumer technology, and even public health policies. According to the framework proposed in Diffusion of Innovations, adoption unfolds through the stages of knowledge, persuasion, decision, implementation, and confirmation, each influenced by social context and communication structure. Furthermore, interactions among the numerous agents involved in disseminating an idea introduce complex dynamics that simpler analytical models cannot fully capture. Consequently, studying innovation diffusion is both challenging and intellectually enriching. The impact of innovation in the healthcare industry is especially important given the life changing impact that it has on humans. 
\section{Background}
The purpose of this study is to investigate how changing the social system affects the diffusion of innovation

\section{Methods}
Good Seth section perhaps
\cite{borracci_agent-based_2018}\cite{orr_diffusion_2013}
To study how the changing of social systems effects the spread of innovation, we designed two policy options for the network creation, a government network, and a local network.
\subsection{Agent Design}

\subsection{Model Design}
\section{Results}
The five adopter categories are:
(1) innovators (2.5) (2) early adopters (13.5), (3) early majority (34), (4) late majority,(34)
and (5) laggards.(16) (p.133)
\section{Conclusion}
\printbibliography
\end{document}
